\documentclass{article}
\usepackage[margin=1in]{geometry}
\setlength{\tabcolsep}{4pt}
\newcommand{\mycite}[1]{\cite{#1}}
\begin{document}

\title{Peer Feedback and Exam Scores}
\maketitle

\begin{abstract}
We study weekly peer feedback in a first year statistics course with 120 participants.
Students who joined the program reached an average of 87.4\% on the final exam and finished 12\% faster.
\end{abstract}

\section{Introduction}
Peer feedback is widely used in large courses \citep[see][p.~3]{good2021, oldyear}.
Prior work reports mixed effects \citet{wrongtitle2020} and \parencite{nolink2017}.
Some studies find no effect at all \cite{fake2023,
  kim2022}.
% An older survey \cite{commented2020} is left out.
About 85\% of instructors use rubrics \cite{missing2021}.
A related design appears in \cite{smith2019a} and \cite{smith2019b}.
See also \cite{kim2022b, chen2020, noyear, badoi2020, subtitled2019}.

\section{Method}
We recruited participants ($n = 120$) from two sections.
Of these, 45\% were first year students, 30\% were second year students and 20\% were repeating the course.
Figure 2 shows the timeline and Table~\ref{tab:main} lists the outcomes.

\section{Results}
Of the remaining participants ($n = 112$), the feedback group scored 85.2\% on the final exam.
Completion time dropped by 12.3\% compared with the previous term.

\begin{table}[t]
\caption{Final exam outcomes for 3 groups.}
\label{tab:main}
\begin{tabular}{lcc}
\toprule
Group & Score (\%) & Time (min) \\
\midrule
Feedback & 85.2 & 41.5 \\
Control & 79.8 & 47.3 \\
\bottomrule
\end{tabular}
\end{table}

\bibliographystyle{plain}
\bibliography{refs}
\end{document}
